% !TEX root = ../main.tex

\chapter{人工智能工具使用情况说明}

\newcolumntype{C}[1]{>{\centering\arraybackslash}p{#1}}
\newcolumntype{L}[1]{>{\raggedright\arraybackslash}p{#1}}
\newcommand{\appendixblank}[1]{\underline{\makebox[#1][l]{}}}

\noindent
\textbf{毕业设计（论文）标题：基于动作扩散策略的
四足机器人视觉目标导航}

\vspace{0.8em}
\noindent
\textbf{学生姓名：陈样}
\hfill
\textbf{学号：522031910056}
\hfill
\textbf{指导教师：雷华明}

\vspace{2em}
\noindent{\zihao{3}\bfseries 一、本文使用的 AI 工具清单}

\vspace{0.8em}
\begingroup
    \centering
    \renewcommand{\arraystretch}{1.65}
    \setlength{\tabcolsep}{4pt}
    \setlength{\LTleft}{0pt}
    \setlength{\LTright}{0pt}
    \begin{longtable}{|C{0.09\linewidth}|C{0.15\linewidth}|C{0.10\linewidth}|C{0.20\linewidth}|L{0.32\linewidth}|}
        \hline
        \textbf{序号} &
        \textbf{AI工具名称} &
        \textbf{版本/型号} &
        \textbf{毕业设计（论文）中应用部分} &
        \centering\arraybackslash\textbf{人为判断过程说明} \\
        \hline
        \endfirsthead
        \hline
        \textbf{序号} &
        \textbf{AI工具名称} &
        \textbf{版本/型号} &
        \textbf{毕业设计（论文）中应用部分} &
        \centering\arraybackslash\textbf{人为判断过程说明} \\
        \hline
        \endhead
        1 & DeepSeek & V4 & C2 & 人工对实验目标把控，调用 AI 工具进行本地实验与调优建议\\
        \hline
        2 & DeepSeek & V4 & D1 & 人工把握论文大纲与逻辑，AI辅助\\
        \hline
        3 & DeepSeek & V4 & D2 & 人工核对正文内容与实验结果的一致性，AI仅用于辅助表述整理。 \\
        \hline
        4 & DeepSeek & V4 & D3 & 人工逐段检查语义、术语和格式，AI用于语言润色建议。 \\
        \hline
        5 & DeepSeek & V4 & E1 & 人工核验参考文献来源与格式，AI用于辅助检查条目完整性。 \\
        \hline
        6 & DeepSeek & V4 & E2 & 人工确认图表内容与论文数据对应关系，AI用于辅助生成和排版建议。 \\
        \hline
    \end{longtable}
\endgroup

\vspace{1.2em}
\noindent\textbf{注释：}

\vspace{0.4em}
\noindent
\textsuperscript{a}可以根据实际情况增删行数；

\vspace{0.4em}
\noindent
\textsuperscript{b}可以对照“应用部分对照表”填写对应的编号表，若超过编号表内容，请根据实际填写；

\vspace{0.4em}
\noindent
\textsuperscript{c}人为判断过程指的是依据何种方式判断 AI 工具生成的内容的准确性。

\vspace{2em}
\noindent{\zihao{3}\bfseries 毕业设计（论文）中应用部分对照表}

\vspace{0.8em}
\begin{center}
    \centering
    \renewcommand{\arraystretch}{1.5}
    \setlength{\tabcolsep}{8pt}
    \begin{tabular}{|L{0.20\linewidth}|L{0.66\linewidth}|}
        \hline
        A.研究准备阶段 & A1.选题与方向设定；A2.文献检索与筛选 \\
        \hline
        B.研究设计阶段 & B1.研究思路设计；B2.方法论构建；B3.研究工具开发 \\
        \hline
        C.研究实施阶段 & C1.实验/数据分析；C2.代码编写与调试；C3.数学建模与计算 \\
        \hline
        D.论文写作阶段 & D1.大纲与逻辑优化；D2.论文正文写作；D3.语言与语法润色 \\
        \hline
        E.其他辅助环节 & E1.参考文献与格式；E2.图表/多媒体生成；E3.学术反馈与模拟答辩 \\
        \hline
    \end{tabular}
\end{center}

\vspace{2.5em}
\begin{center}
    {\zihao{3}\bfseries 学术诚信声明}
\end{center}

\vspace{1em}
\noindent
本人承诺：论文核心观点、实验数据及结论均独立完成，AI 工具仅用于辅助性工作，已明确说明所有生成内容并在文中对应位置进行脚注标注。

\vspace{3em}
\begin{flushright}
    承诺人：陈样
    日期：2026年5月7日
\end{flushright}
